
\section*{Discussion}

In this study, we evaluated the predictive performance of two pressure-based biodiversity models using a global dataset. Our results highlight limits to spatially explicit predictions across models (\autoref{fig:summary_results}), related to data availability and model design. These findings suggest that the generality of large-scale biodiversity models remains a major challenge\cite{spake2022}. Some generalization could be attained within contexts where there was enough data (\autoref{fig:bhm_alpha_groups}), but transferability to new contexts remains elusive (\autoref{fig:bhm_alpha_distr_shifts}). This aligns with prior results at smaller spatial and taxonomic scales\cite{roberts2017,meyer2021,meyer2022}. Given the demand for global biodiversity indicators\cite{burgess2024}, the systematic assessment presented here fills an important knowledge gap. Our results have implications for the use and interpretation of pressure-based models and indicators. These insights can support data collection and model development efforts to further strengthen the critical role of biodiversity models for global policy and decision-making \cite{zurell2025, purvis2025}.

By comparing two structurally different models, we are able to shed light on different aspects of the model generality challenge. Lingering data gaps and biases on a global level \cite{hughes2021,chapman2024} puts fundamental constraints on large-scale biodiversity models. This challenge is clearly shown by the decrease in predictive performance of the BTMs between interpolation and extrapolation (\figpanel{fig:summary_results}{a,b} and \figpanel{fig:bhm_alpha_distr_shifts}{b,c}). The data distribution shifts (\figpanel{fig:bhm_alpha_distr_shifts}{d,e}) that forced the models to make out-of-distribution predictions cannot be adequately addressed without additional data. More representative data implies fewer places where extrapolation is required, and would also enable interpolation within increasingly relevant contexts. To illustrate the second point, consider that the parameters in the BTMs were at the level of 'birds in moist broadleaf forests of the Neotropics', which implies averaging over a great deal of underlying differences in distributions and responses to pressures. Indeed, we saw large variation in performance between hierarchical groups (\figpanel{fig:bhm_alpha_groups}{a}) and across countries (\autoref{fig:bhm_country_maps}), related to data availability and heterogeneity.

The results also highlight the role of model design choices. The SBMs had consistently low performance across interpolation and extrapolation (\autoref{fig:summary_results} and \figpanel{fig:lmm_alpha_deepdive}{b,c}). The use of study-level random effects is sensible for estimating average effects of anthropogenic drivers \cite{harrison2018,pinheiro2000}. However, they also prevent generalization, since models attribute signal to variables that are discarded at the prediction stage (\figpanel{fig:lmm_alpha_deepdive}{d}). This underscores that explanatory power is not necessarily an indicator of predictive power\cite{shmueli2010,tredennick2021}. In comparison, the structure and covariates of the BTMs enabled some, albeit limited, generalization (\FigsTwo{fig:summary_results}{}{fig:bhm_alpha_groups}{a,b}). It is not surprising that a model trained using more contextually explicit information is better at predicting relative levels of alpha diversity (\figpanel{fig:summary_results}{a}). But even if the main goal was to differentiate between sites with different pressures within a given ecological context, our results indicate that the BTMs had an advantage due to their varying-group parameters (\figpanel{fig:summary_results}{b,c,d}). All of this naturally hinges on data being reasonable representative in the biogeographic-taxonomic groups modeled. This is unlikely to be the case everywhere, especially in groups with fewer observations. Still, this hierarchical structure provides a dynamic mechanism for capturing the benefits of increasing data across contexts. Although something similar could be achieved by training separate SBMs per group, it would be less dynamic and not leverage the full data as effectively. For spatial prediction, we argue that it is generally preferable to explicitly model both the natural and anthropogenic factors at play. It should also be noted that our goal was to give an illustrative comparison between the two modeling approaches, not exhaustively optimize performance; that would have involved using more extensive biodiversity, environmental, and anthropogenic data.

Scientifically robust indicators are essential for achieving the goals of the GBF \cite{affinito2024,gbfmf2025}. Global implementations are appealing because they are readily applicable and offer promises of consistency\cite{ipbes2019}. Our results corroborate that, on average, land-use change is a strong driver of biodiversity loss (\figpanel{fig:lmm_alpha_deepdive}{e}), in line with many previous analyses \cite{keck2025,jaureguiberry2022,newbold2015,depalma2021,bevan2024}. Yet, the low model accuracies emphasize the challenging gap between effect size inference and spatially explicit predictions. Maps where spatial resolution is not matched by predictive power can incur concrete risks, such as suboptimal conservation priorities or inaccurate sustainability reports. To effectively support decision-making, the limitations of model-based indicators must be transparent to users. Countries with extensive data can overcome certain limitations by implementing their own models and indicators\cite{gbfmf2025,affinito2024}. However, that would still leave data-poor regions with a lack of accurate biodiversity insights on local scales\cite{clements2025}. Here, resources must be allocated to scale up collection of standardized biodiversity data in prioritized areas and taxonomic groups\cite{hughes2021,gonzalez2023}, leveraging cost-effective technologies like environmental DNA, bioacoustics and camera traps. Similarly, companies that want to ensure accurate decision-making and reporting will often need to collect proprietary data to complement existing tools\cite{goodsell2025}.

%One motivation behind this study was the use of pressure-based models for spatial projections\cite{newbold2015,schipper2020} despite lack of systematic testing\cite{martin2019}.

%In other words, the gap between directional signals and precise estimates can be large.

Our study highlights several research priorities for improving large-scale, predictive biodiversity models. There are clear opportunities to build dynamic and scalable data pipelines that combine large quantities of sampling event data from repositories like the Global Biodiversity Information Facilitaty (GBIF), with the latest remote sensing data from platforms like Google Earth Engine. Contributions of structured sampling data to central repositories should be prioritized, where large-scale datasets are of particular value. More effort should be invested into methodological development of top-down community models, like the ones used in this study, as a complement to species distribution models (SDMs). By aggregating biodiversity data first and then training models, they offer necessary scalability for global applications. Methods that tackle data-sparse contexts and explore the limits of extrapolation are other important avenues of research. Finally, evaluating biodiversity models with such a large scope is a very complex task, and limitations to site-level evaluation and model intercomparison have been highlighted already. This points to the need for evaluation frameworks that can handle scale in multiple dimensions, to assess whether a large-scale model is 'good enough' in a certain context, as opposed to more locally adapted models.

%In conclusion, models are essential for biodiversity monitoring, but our findings demonstrate that the predictive power of pressure-based models has clear limitations. Users need to be aware of these limitations, and the scientific community should prioritize the mobilization of data and development of refined models, to ensure that large-scale biodiversity indicators can reliably guide conservation efforts.
